\begin{pastexam}[label={px:ncku-114-cce-ds-13}]{成大電通所 114 學年度 單選題第 13 題}
（單選；每題 5 分，答錯倒扣 2 分）

Please indicate the algorithm that cannot allow negative weighted edges in finding the shortest path.\\
a.) Jonson's algorithm,\quad b.) Bellman and Ford's Algorithm,\quad c) Floyd Warshall's Algorithm,\quad
d) Dijkstra's Algorithm,\quad e) None of above.
\end{pastexam}

\begin{pastsoln}
答案：d) Dijkstra's Algorithm。

\begin{itemize}
\item d) Dijkstra 每次把目前距離最小的頂點定案，之後不再更新。 這要靠「邊權非負，所以後面的路只會更長」；有負邊時，晚一點才走到的
負邊可能讓已定案的頂點變得更近，答案就錯（資結書第~25.2.5~節）。 例如 $s \to a$ 權重 2、$s \to b$ 權重 3、$b \to a$ 權重 $-2$：
Dijkstra 先定案 $a$（距離 2），但真正的最短距離是 $s \to b \to a = 1$。
\item a) Johnson（原題拼成 Jonson）：先用 Bellman--Ford 算出位能 $h$，把每條邊改成 $w(u,v) + h(u) - h(v) \ge 0$，再對每個頂點跑
Dijkstra。 它就是為了處理負邊而設計的，只要沒有負環（資結書第~25.7~節）。
\item b) Bellman--Ford：允許負邊，還能偵測從起點走得到的負環（資結書第~25.3~節）。
\item c) Floyd--Warshall：允許負邊，只要沒有負環（資結書第~25.6~節）。
\end{itemize}
\end{pastsoln}
